\clearpage
\begin{summarybox}
\textbf{一句话记忆：\quad 不用猜函数长相，只管连续走步，走到函数值全部还原为止。}

\textbf{考场速查：}
\[
\begin{array}{rcl}
 f(x+a)=f(x)&\Rightarrow&T=|a|;\\[3pt]
 f(x+a)=-f(x)&\Rightarrow&T=2|a|;\\[3pt]
 f(x+a)=\dfrac{\pm1}{f(x)}&\Rightarrow&T=2|a|;\\[7pt]
 f(x+a)=f(x+b)&\Rightarrow&T=|a-b|;\\[3pt]
 f(x+a)=c-f(x)&\Rightarrow&T=2|a|;\\[3pt]
 f(x+a)=\dfrac{c}{f(x)}&\Rightarrow&T=2|a|.
\end{array}
\]
每个 $T$ 都是\textbf{一个正周期}，不保证最小。
在 $\pm$ 型中正负号要固定；$c$ 为固定常数，倒数时 $c\ne0$、$f(x)\ne0$；
差值型要有 $a\ne b$；还应检查所有涉及的输入都在定义域中。

\textbf{两种分式也要会识别：}
\[
 \frac{f(x)+1}{f(x)-1}\quad\text{与}\quad\frac{1-f(x)}{1+f(x)}.
\]
若它们分别等于 $f(x+a)$，且每次分母合法、关系对全部允许的位置成立，
则都可推出 $\boxed{f(x+2a)=f(x)}$，因此 $\boxed{T=2|a|}$。

\textbf{四个关键动作：}
\begin{enumerate}[leftmargin=*,itemsep=0.35em]
\item 先判断“往右走一步”后，数值究竟做了什么变化。
\item 若一步没有还原，就把\textbf{同一条等式再用一次}。
\item 若两边都是 $f(x+\text{常数})$，把输入的\textbf{差}算出来。
\item 求 $f(2026)$ 这类值时，用已经证明的周期把大数化到已知的小位置。
\end{enumerate}

\textbf{一个反直觉的提醒：} $f(x+1)=-f(x)$ 推出 $2$ 是周期，却可能还有更小周期；
例如 $\sin(3\pi x)$ 的最小正周期是 $2/3$。
\textbf{与图像对称性的联系：}两条不同对称轴 $x=p$、$x=q$ 可推出一个周期 $2|p-q|$。
\end{summarybox}
